\documentclass[a4paper, 12pt, twoside]{article}
\usepackage[utf8]{inputenc}
\usepackage[T1]{fontenc}
\usepackage[francais]{babel}
\usepackage{lmodern}
\usepackage{ae,aecompl}
\usepackage[top=2.5cm, bottom=2cm,
			left=3cm, right=2.5cm,
			headheight=15pt]{geometry}
\usepackage{graphicx}
\usepackage{eso-pic}
\usepackage{array}
\usepackage{hyperref}
\usepackage{float}
\usepackage{listings}
\usepackage{xcolor}
\usepackage{amsmath}   % \mathbb, \to, etc.
\usepackage{amssymb}   % \mathbb{R}

\lstset{
    basicstyle=\ttfamily\small,
    breaklines=true,
    frame=single,
    language=Python,
    keywordstyle=\color{blue},
    commentstyle=\color{green!50!black},
    stringstyle=\color{red},
    numbers=left,
    numberstyle=\tiny,
    numbersep=5pt
}

\input{pagedegarde}
\title{Analyse du réseau de bus en Île-de-France}
\membrea{Rayane MEKSEM}
\membreb{Ryan JARDIN}
\membrec{Monju DEY}

\begin{document}
\pagedegarde

% ----------------------------------------------------------------
\section*{Remerciements}
Nous tenons à remercier notre enseignant responsable du module Graphes et OpenData pour son accompagnement tout au long de ce projet. Nous remercions également Île-de-France Mobilités pour la mise à disposition des données ouvertes relatives au réseau de transport en commun francilien, sans lesquelles ce travail n'aurait pas été possible.

\newpage
\tableofcontents
\newpage

% ----------------------------------------------------------------
\section{Introduction}

\paragraph{Attention !} Vous devez impérativement utiliser ce template pour réaliser votre rapport de projet d'harmonisation. Vous n'avez pas le droit de changer la taille, la police, l'interligne ou les marges.\\

L'Île-de-France est la région la plus peuplée de France avec près de 12 millions d'habitants répartis sur 8 départements. Son réseau de transport en commun, géré par Île-de-France Mobilités, est l'un des plus denses d'Europe. Il comprend notamment un vaste réseau de bus composé de plus de 1\,500 lignes et 25\,000 arrêts.

Ce projet exploite les données ouvertes publiées par Île-de-France Mobilités \cite{IDFM} au format GTFS (\textit{General Transit Feed Specification}) \cite{GTFSDOC} pour modéliser le réseau de bus francilien sous forme de graphe et en analyser les propriétés structurelles à l'aide d'algorithmes classiques.

Le Tableau~\ref{tab:gtfs} présente les fichiers GTFS utilisés. La Figure~\ref{fig:subgraph} donne un aperçu du graphe construit.

\begin{table}[h]
\begin{center}
\begin{tabular}{|l|l|r|}
\hline
\textbf{Fichier} & \textbf{Contenu} & \textbf{Entrées (approx.)} \\
\hline
\texttt{stops.txt}       & Arrêts de bus           & $\sim$25\,000     \\
\texttt{routes.txt}      & Lignes de bus           & $\sim$1\,500      \\
\texttt{trips.txt}       & Trajets planifiés       & $\sim$500\,000    \\
\texttt{stop\_times.txt} & Horaires par arrêt      & $\sim$8\,000\,000 \\
\hline
\end{tabular}
\end{center}
\caption{Fichiers GTFS d'Île-de-France Mobilités utilisés pour le projet}
\label{tab:gtfs}
\end{table}

Et ne jamais oublier que lorsqu'on place un tableau, on doit l'utiliser comme ici avec le tableau Table~\ref{tab:gtfs}. Et on fait la même chose avec une image, comme avec la Figure~\ref{fig:subgraph}.

\begin{figure}[H]
\centering
\includegraphics[width=0.85\textwidth]{figures/fig1_subgraph.pdf}
\caption{Sous-graphe des 150 arrêts les plus connectés du réseau de bus IDF. La couleur et la taille des nœuds sont proportionnelles à leur degré.}
\label{fig:subgraph}
\end{figure}

% ----------------------------------------------------------------
\section{Environnement de travail}

Le projet a été développé dans l'environnement suivant :

\begin{itemize}
    \item \textbf{Système d'exploitation :} Ubuntu 20.04 LTS 
    \item \textbf{Langage de programmation :} Python 3.9
    \item \textbf{Environnement de développement :} Visual Studio Code, Jupyter Notebook
    \item \textbf{Gestion de version :} GitHub 
    \item \textbf{Matériel :} Ordinateurs portables personnels (8 Go RAM, processeur Intel Core i5)
\end{itemize}

La collaboration a été facilitée par GitHub pour le partage du code et Discord pour la communication.

% ----------------------------------------------------------------
\section{Description du projet et objectifs}

	\subsection{Modélisation du réseau sous forme de graphe}

	Le réseau de bus est modélisé comme un \textbf{graphe orienté pondéré} $G = (V, E, w)$ où :
	\begin{itemize}
	    \item $V$ est l'ensemble des sommets : chaque arrêt de bus est un noeud.
	    \item $E$ est l'ensemble des arêtes orientées : deux arrêts $u$ et $v$ sont reliés par un arc $u \rightarrow v$ s'il existe un trajet de bus qui les dessert consécutivement.
	    \item $w : E \rightarrow \mathbb{R}^+$ est la fonction de poids : le poids d'un arc est le temps de trajet moyen (en secondes) entre les deux arrêts, calculé à partir du fichier \texttt{stop\_times.txt}.
	\end{itemize}

	Le graphe est \textbf{orienté} car les lignes de bus ne sont pas nécessairement symétriques (sens aller $\neq$ sens retour) et \textbf{pondéré} car les temps de trajet varient d'un tronçon à l'autre.

	\subsection{Objectifs du projet}

	Les objectifs concrets sont les suivants :
	\begin{itemize}
	    \item Charger et parser les données GTFS d'Île-de-France Mobilités.
	    \item Construire le graphe du réseau de bus avec la bibliothèque NetworkX \cite{NETWORKX}.
	    \item Calculer des indicateurs structurels : degré moyen, distribution des degrés, connexité.
	    \item Identifier les arrêts les plus importants par la centralité d'intermédiarité (\textit{betweenness centrality}).
	    \item Trouver le plus court chemin entre deux arrêts avec l'algorithme de Dijkstra \cite{DIJKSTRA}.
	    \item Visualiser et interpréter les résultats sous forme de figures.
	\end{itemize}

% ----------------------------------------------------------------
\section{Bibliothèques, Outils et technologies}

\begin{table}[h]
\begin{center}
\begin{tabular}{|l|c|l|}
\hline
\textbf{Bibliothèque} & \textbf{Version} & \textbf{Usage} \\
\hline
\texttt{NetworkX}   & 2.6  & Construction et analyse du graphe \\
\texttt{Pandas}     & 1.3  & Chargement et filtrage des fichiers GTFS \\
\texttt{Matplotlib} & 3.4  & Génération des figures \\
\texttt{Folium}     & 0.12 & Cartographie interactive \\
\texttt{NumPy}      & 1.21 & Calculs numériques \\
\hline
\end{tabular}
\end{center}
\caption{Bibliothèques Python utilisées dans le projet}
\label{tab:libs}
\end{table}

Les données sont au format GTFS, un ensemble de fichiers CSV standardisés (voir Table~\ref{tab:gtfs}). Pandas permet de les charger et filtrer efficacement avant de construire le graphe avec NetworkX.

% ----------------------------------------------------------------
\section{Travail réalisé}

La liste des fonctionnalités prévues et l'état de réalisation sont présentés dans le Tableau~\ref{tab:fonc}.

\begin{table}[h]
\begin{center}
\begin{tabular}{|p{7cm}|c|p{4.5cm}|}
\hline
\textbf{Fonctionnalité} & \textbf{Réalisée} & \textbf{Commentaire} \\
\hline
Chargement des données GTFS             & Oui & Filtrage sur bus (type 3) \\
Construction du graphe orienté pondéré  & Oui & NetworkX DiGraph \\
Visualisation du sous-graphe            & Oui & Figure \ref{fig:subgraph} \\
Distribution des degrés                 & Oui & Figure \ref{fig:degres} \\
Centralité (betweenness)                & Oui & Figure \ref{fig:central} \\
Plus court chemin (Dijkstra)            & Oui & Figure \ref{fig:dijkstra} \\
Composantes connexes                    & Oui & Figure \ref{fig:compo} \\
Interface graphique utilisateur         & Non & Manque de temps \\
\hline
\end{tabular}
\end{center}
\caption{Liste des fonctionnalités prévues et état de réalisation}
\label{tab:fonc}
\end{table}

La fonctionnalité d'interface graphique n'a pas été réalisée faute de temps : la gestion du volume des données GTFS et le débogage des algorithmes ont été plus chronophages que prévu.\\

\noindent\textbf{Résultats obtenus.} Après construction à partir des données GTFS, le graphe présente les caractéristiques structurelles suivantes :

\begin{table}[h]
\begin{center}
\begin{tabular}{|l|r|}
\hline
\textbf{Indicateur} & \textbf{Valeur} \\
\hline
Nombre de sommets (arrêts)                & $\approx 24\,800$ \\
Nombre d'arêtes (liaisons)                & $\approx 62\,000$ \\
Degré moyen                               & $\approx 5.0$     \\
Nombre de composantes faiblement connexes & 312               \\
Taille de la composante principale        & 98\% des sommets  \\
\hline
\end{tabular}
\end{center}
\caption{Caractéristiques structurelles du graphe du réseau de bus IDF}
\label{tab:stats}
\end{table}

\noindent\textbf{Figure 2 — Distribution des degrés.}\\
La Figure~\ref{fig:degres} montre que la majorité des arrêts possède un faible degré (entre 2 et 6), tandis qu'un petit nombre de noeuds présentent un degré très élevé. Cette distribution hétérogène est caractéristique des réseaux de transport organisés autour de hubs.

\begin{figure}[H]
\centering
\includegraphics[width=\textwidth]{figures/fig2_degree_distribution.pdf}
\caption{Distribution des degrés du réseau. À gauche : histogramme (ligne rouge = degré moyen). À droite : même distribution en échelle log-log.}
\label{fig:degres}
\end{figure}

\noindent\textbf{Figure 3 — Centralité d'intermédiarité.}\\
La Figure~\ref{fig:central} présente le top 15 des arrêts les plus centraux. Ces arrêts correspondent aux grands hubs de correspondance ; leur suppression fragiliserait fortement la connexité du réseau.

\begin{figure}[H]
\centering
\includegraphics[width=0.85\textwidth]{figures/fig3_centrality.pdf}
\caption{Top 15 des arrêts selon la centralité d'intermédiarité (\textit{betweenness centrality}), calculée par échantillonnage ($k = 500$).}
\label{fig:central}
\end{figure}

\noindent\textbf{Figure 4 — Plus court chemin par l'algorithme de Dijkstra.}\\
L'algorithme de Dijkstra \cite{DIJKSTRA} est appliqué sur le graphe pondéré (poids = temps en secondes). La Figure~\ref{fig:dijkstra} illustre le chemin optimal calculé entre deux arrêts.

\begin{figure}[H]
\centering
\includegraphics[width=0.85\textwidth]{figures/fig4_shortest_path.pdf}
\caption{Plus court chemin (en rouge) entre deux deux arrêts au hasard calculé par Dijkstra. Noeuds rouges : départ/arrivée. Noeuds oranges : arrêts intermédiaires.}
\label{fig:dijkstra}
\end{figure}

\noindent\textbf{Figure 5 — Composantes connexes.}\\
La Figure~\ref{fig:compo} montre que le réseau est quasi-connexe : la composante géante englobe $\approx$98\% des arrêts. Les petites composantes isolées correspondent à des arrêts en zones périphériques sans correspondance avec le reste du réseau.

\begin{figure}[H]
\centering
\includegraphics[width=\textwidth]{figures/fig5_components.pdf}
\caption{Analyse des composantes connexes. À gauche : répartition (la composante géante représente 98\% des noeuds). À droite : distribution des tailles en échelle logarithmique.}
\label{fig:compo}
\end{figure}

\noindent\textbf{Extrait de code — construction du graphe :}

\begin{lstlisting}[caption={Chargement GTFS et construction du graphe NetworkX}]
import pandas as pd
import networkx as nx

stops      = pd.read_csv("data/stops.txt", dtype=str)
stop_times = pd.read_csv("data/stop_times.txt", dtype=str)
stop_times["stop_sequence"] = stop_times["stop_sequence"].astype(int)

G = nx.DiGraph()
for _, row in stops.iterrows():
    G.add_node(row["stop_id"], name=row["stop_name"],
               lat=float(row["stop_lat"]),
               lon=float(row["stop_lon"]))

st = stop_times.sort_values(["trip_id", "stop_sequence"])
for trip_id, group in st.groupby("trip_id"):
    group = group.reset_index(drop=True)
    for i in range(len(group) - 1):
        u = group.loc[i,   "stop_id"]
        v = group.loc[i+1, "stop_id"]
        w = parse_time(group.loc[i+1, "arrival_time"]) \
          - parse_time(group.loc[i,   "departure_time"])
        if w > 0:
            G.add_edge(u, v, weight=w)
\end{lstlisting}

\noindent\textbf{Répartition réelle du travail entre les membres du groupe :}

\begin{table}[h]
\begin{center}
\begin{tabular}{|l|p{9cm}|}
\hline
\textbf{Membre} & \textbf{Tâches réalisées} \\
\hline
MEKSEM & Chargement GTFS, construction du graphe, calcul des indicateurs, figures 1 et 2 \\
JARDIN & Dijkstra, betweenness centrality, figures 3/4/5 \\
DEY & rédaction du rapport \\
\hline
\end{tabular}
\end{center}
\caption{Répartition réelle du travail entre les membres du groupe}
\label{tab:travail}
\end{table}

% ----------------------------------------------------------------
\section{Difficultés rencontrées}

\begin{enumerate}
    \item \textbf{Volume des données :} Le fichier \texttt{stop\_times.txt} contient plusieurs millions de lignes. Le chargement direct saturait la mémoire vive. Nous avons résolu ce problème en chargeant les données par blocs (\textit{chunks}) avec Pandas et en filtrant immédiatement les lignes de bus (type GTFS = 3).

    \item \textbf{Heures supérieures à 24h :} Le format GTFS autorise des horaires comme \texttt{25:30:00} pour les services nocturnes qui chevauchent minuit. Notre fonction de conversion a été adaptée pour gérer ces cas.

    \item \textbf{Temps de calcul de la betweenness centrality :} Sur un graphe de 25\,000 noeuds, le calcul exact est prohibitif ($O(n \cdot m)$). Nous avons utilisé l'approximation par échantillonnage de NetworkX ($k = 500$ noeuds sources), réduisant le temps de calcul de plusieurs heures à quelques minutes.

    \item \textbf{Visualisation du graphe complet :} Afficher 25\,000 noeuds avec Matplotlib produisait une image illisible. Nous avons choisi de représenter uniquement les 150 arrêts les plus connectés (Figure~\ref{fig:subgraph}).
\end{enumerate}

% ----------------------------------------------------------------
\section{Bilan}

	\subsection{Conclusion}

	Ce projet nous a permis de mettre en pratique la théorie des graphes sur un jeu de données réel et volumineux. Nous avons modélisé le réseau de bus d'Île-de-France sous forme de graphe orienté pondéré, calculé ses indicateurs structurels et identifié les arrêts les plus stratégiques.

	Les cinq figures produites montrent que le réseau présente les caractéristiques typiques d'un réseau de transport : une composante géante, une distribution des degrés hétérogène, et quelques hubs très centraux dont dépend fortement la cohérence globale du réseau.

	\subsection{Perspectives}

	\begin{itemize}
	    \item \textbf{Multi-modal :} Intégrer le métro, le RER et le Transilien pour construire un graphe de transport complet.
	    \item \textbf{Temporalité :} Distinguer les heures de pointe et les heures creuses dans les poids des arêtes pour des itinéraires plus réalistes.
	    \item \textbf{Interface web :} Développer une application Flask permettant à l'utilisateur de calculer un itinéraire de façon interactive.
	    \item \textbf{Résilience :} Simuler la suppression de noeuds centraux pour évaluer la robustesse du réseau face à des pannes ou perturbations.
	\end{itemize}

\newpage
% ----------------------------------------------------------------
\section{Bibliographie}
\renewcommand{\bibname}{}
\renewcommand{\refname}{}
\begin{thebibliography}{4}
   \bibitem[LAM94]{lam1} L. LAMPORT, {\it \LaTeX : A Document preparation system}, Addison-Wesley, 1994.
   \bibitem[HAG08]{NETWORKX} A. HAGBERG, P. SWART, D. SCHULT, {\it Exploring Network Structure, Dynamics, and Function using NetworkX}, Proceedings of the 7th Python in Science Conference, 2008.
   \bibitem[DIJ59]{DIJKSTRA} E.W. DIJKSTRA, {\it A note on two problems in connexion with graphs}, Numerische Mathematik, vol. 1, pp. 269--271, 1959.
   \bibitem[GGL06]{GTFS} Google LLC, {\it GTFS Reference -- General Transit Feed Specification}, 2006.
\end{thebibliography}

\newpage
% ----------------------------------------------------------------
\section{Webographie}
\begin{thebibliography}{4}
   \bibitem[IDFM]{IDFM} \url{https://data.iledefrance-mobilites.fr} --- Portail OpenData d'Île-de-France Mobilités
   \bibitem[NXD]{NETWORKXDOC} \url{https://networkx.org/documentation/stable/} --- Documentation NetworkX
   \bibitem[GTFD]{GTFSDOC} \url{https://gtfs.org/reference/static/} --- Spécification GTFS
   \bibitem[FOL]{FOLIUMDOC} \url{https://python-visualization.github.io/folium/} --- Documentation Folium
\end{thebibliography}

\newpage
% ----------------------------------------------------------------
\section{Annexes}
\appendix
\makeatletter
\def\@seccntformat#1{Annexe~\csname the#1\endcsname:\quad}
\makeatother

	\section{Cahier des charges}

	\textbf{Objectif général :} Analyser le réseau de bus en Île-de-France à l'aide de la théorie des graphes et des données ouvertes GTFS (licence ODbL).\\

	\noindent\textbf{Données :} Fichiers GTFS publiés par Île-de-France Mobilités, téléchargeables sur \url{https://data.iledefrance-mobilites.fr}.\\

	\noindent\textbf{Livrables attendus :}
	\begin{itemize}
	    \item Un programme Python documenté (\texttt{generate\_graphs.py}).
	    \item Le présent rapport au format PDF.
	    \item Cinq figures générées automatiquement (dossier \texttt{figures/}).
	\end{itemize}

	\section{Exemple d'exécution du projet}

\begin{lstlisting}[language=bash, caption={Exécution du script principal}]
$ pip install networkx pandas matplotlib folium numpy
$ python generate_graphs.py

Chargement des donnees GTFS...
  24812 arrets | 1487 lignes bus | 7923045 horaires
Construction du graphe : 24812 noeuds, 61934 aretes
Figure 1 : sous-graphe...          -> figures/fig1_subgraph.pdf
Figure 2 : distribution degres...  -> figures/fig2_degree_distribution.pdf
Figure 3 : centralite...           -> figures/fig3_centrality.pdf
Figure 4 : plus court chemin...    -> figures/fig4_shortest_path.pdf
Figure 5 : composantes connexes... -> figures/fig5_components.pdf
Toutes les figures sont dans ./figures/
\end{lstlisting}

	\section{Manuel utilisateur}

	\textbf{Prérequis :} Python 3.8 ou supérieur, pip.\\

	\noindent\textbf{Installation des dépendances :}
\begin{lstlisting}[language=bash]
pip install networkx pandas matplotlib folium numpy
\end{lstlisting}

	\noindent\textbf{Téléchargement des données :}
	\begin{enumerate}
	    \item Aller sur \url{https://data.iledefrance-mobilites.fr}
	    \item Rechercher « GTFS » et télécharger \texttt{IDFM-gtfs.zip}
	    \item Dézipper dans le dossier \texttt{data/} du projet
	\end{enumerate}

	\noindent\textbf{Lancement :}
\begin{lstlisting}[language=bash]
python3 generate_graphs.py
\end{lstlisting}
	Les figures PDF sont générées dans \texttt{figures/} et incluses automatiquement dans le rapport lors de la compilation avec \texttt{pdflatex}.

\end{document}
